% Solution to Problem 1 — (Lack of) quasi-shift-invariance of the Phi^4_3 measure (Hairer)
% v8, corrected against First Proof paper (Hairer-Peroni)
% Corrects a sign-level error: the true answer is MUTUAL SINGULARITY (NO), not equivalence.

\documentclass[11pt]{article}
\usepackage[margin=1in]{geometry}
\usepackage{amsmath, amssymb, amsthm}
\usepackage{mathtools}
\usepackage{xcolor}
\usepackage{hyperref}
\hypersetup{colorlinks=true,
            linkcolor=blue!50!black,
            citecolor=blue!50!black,
            urlcolor=blue!50!black}

\newcommand{\T}{\mathbb{T}}
\newcommand{\R}{\mathbb{R}}
\newcommand{\N}{\mathbb{N}}
\newcommand{\C}{\mathcal{C}}
\newcommand{\D}{\mathcal{D}}
\newcommand{\E}{\mathbb{E}}
\newcommand{\Hp}{H}      % Hermite
\newcommand{\PN}{P_N}
\newcommand{\Lin}{\mathbf{l}}   % stationary linear solution (placeholder for tree symbol)
\newcommand{\Lsq}{\mathbf{q}}   % Wick square
\newcommand{\Lcb}{\mathbf{c}}   % Wick cube
\newcommand{\Iqc}{\mathbf{C}}   % solution driven by the cube

\newtheorem{theorem}{Theorem}
\newtheorem{lemma}{Lemma}
\newtheorem{proposition}{Proposition}
\newtheorem{corollary}{Corollary}

\title{Solution to Problem 1: (Lack of) Quasi-Shift-Invariance\\
of the $\Phi^4_3$ Measure}
\author{}
\date{May 27, 2026}

\begin{document}
\maketitle

\section*{Statement and Answer}

Let $\T^3$ be the unit $3$-torus and $\mu$ the $\Phi^4_3$ measure on
$\D'(\T^3)$. For smooth $\psi\not\equiv0$ let $T_\psi(u)=u+\psi$, and let
$T_\psi^*\mu$ be the pushforward.

\medskip
\noindent\textbf{Theorem.} \emph{For every smooth $\psi\not\equiv0$ and every
admissible choice of mass and (nonzero) coupling constant, the measures $\mu$ and
$T_\psi^*\mu$ are mutually \textbf{singular}.}

\medskip
\noindent The answer is \textbf{NO}: $\mu$ and $T_\psi^*\mu$ are \emph{not}
equivalent. This is a simplified form of a result of Hairer--Peroni and rests on
ideas of Hairer--Kusuoka--Nagoji.

\medskip
\noindent\textbf{Warning on a fatal false premise.} It is tempting to write
$d\mu/d\mu_0=Z^{-1}e^{-\mathcal R(u)}$ against the Gaussian free field $\mu_0$
and deduce quasi-invariance by a Cameron--Martin/Wick-binomial computation. This
is \emph{wrong in three dimensions}: Glimm showed that, unlike $d=2$, the
$\Phi^4_3$ measure $\mu$ and the free field are mutually \emph{singular} in
$d=3$ (the borderline dimension is $8/3$). There is no such Radon--Nikodym
density, so any argument built on it collapses. The correct proof exhibits an
event separating $\mu$ from $T_\psi^*\mu$, driven by a genuinely
$3$-dimensional logarithmic divergence.

\section*{Notation}

Fix a space--time white noise $\xi$ on $\R\times\T^3$ and let $\Lin$ be the
stationary solution of $(\partial_t+1-\Delta)\Lin=\xi$. Write $\PN=P_{\le N}$ for
the projection onto Fourier modes $|k|\le N$, $\Lin_N=\PN\Lin$, and
$c_N=\E\Lin_N^2$. Let $\Hp_n$ be the Hermite polynomials normalized by
$\Hp_0\equiv1$, $\Hp_n'=n\Hp_{n-1}$, $\E\Hp_n(Z,1)=0$ for standard normal $Z$. Set
\[
\Lsq=\lim_{N}\Hp_2(\Lin_N,c_N),\qquad \Lcb=\lim_N\Hp_3(\Lin_N,c_N),
\]
the Wick square and cube (convergent in $\C^{-1-2\kappa}$ and $\C^{-3/2-3\kappa}$
respectively), and let $\Iqc$ solve $(\partial_t+1-\Delta)\Iqc=\Lcb$. We fix
$0<\kappa\le 1/100$. By \cite{HKN24}, $\Lcb$ and $\Iqc$ are a.s.\ continuous in
time with values in $\C^{-1/2-\kappa}$ and $\C^{1/2-3\kappa}$.

We will need the additional constant
\begin{equation}\label{eq:cN2}
c_{N,2}:=\E\bigl[\Lsq_N\,\Iqc_N\bigr],\qquad \Lsq_N:=\PN\Lsq,\ \Iqc_N:=\PN\Iqc.
\end{equation}

\section*{The decomposition of $\mu$}

\begin{proposition}[\cite{HM18,EW24,CC18}]\label{prop:decomp}
There is a stationary process $v$, a.s.\ continuous in time with values in
$\C^{1-2\kappa}(\T^3)$, such that
\begin{equation}\label{eq:u}
u=\Lin-\Iqc+v
\end{equation}
is stationary with fixed-time law $\mu$. Moreover $v$ admits the paracontrolled
form
\begin{equation}\label{eq:v}
v=-3\bigl((v-\Iqc)\prec\Lin\bigr)+v^\sharp,
\end{equation}
where $v^\sharp\in\C^{1+4\kappa}$ and $\prec$ is Bony's paraproduct.
\end{proposition}

\section*{The distinguishing event}

For $\gamma>0$ define
\[
B_\gamma:=\Bigl\{u\in\D':\ \lim_{N\to\infty}\bigl\langle(\log N)^{-\gamma}\bigl(\Hp_3(\PN u;c_N)+9c_{N,2}\PN u\bigr),\psi\bigr\rangle_{\T^3}=0\Bigr\},
\]
where $N$ runs over $2^{\N}$. The next two lemmas (proved as in
\cite[Lems.~3.11--3.12]{HKN24}) supply the analytic input.

\begin{lemma}\label{lem:cube}
For $\gamma>\tfrac12$ and each fixed $t>0$,
$(\log N)^{-\gamma}\,{:}\Lin_N^3{:}(t)\to0$ a.s.\ in $\C^{-3/2}(\T^3)$ and in
$L^p(\Omega;\C^{-3/2})$ for all $p>0$.
\end{lemma}

\begin{proof}[Sketch]
With $\beta=-d/2=-3/2$ and the embedding $W^{\beta,2p}\hookrightarrow
\C^{\beta-d/2p}$, one bounds
$\E\|(\log N)^{-\gamma}{:}\Lin_N^3{:}\|_{W^{-3/2,2p}}^{2p}$ by a constant times
\[
(\log N)^{-2\gamma}\sum_{|\omega_i|\le N}\langle\omega_1+\omega_2+\omega_3\rangle^{-3}\prod_{i}\langle\omega_i\rangle^{-2}
\ \lesssim\ (\log N)^{-2\gamma+1},
\]
which is summable along $N\in2^\N$ once $\gamma>\tfrac12$; Borel--Cantelli
concludes.
\end{proof}

\begin{lemma}\label{lem:logdiv}
For $N$ large, $c_{N,2}\gtrsim\log N$.
\end{lemma}

\begin{proof}[Sketch]
Expanding \eqref{eq:cN2} and integrating the heat kernels in time,
\[
c_{N,2}\gtrsim\sum_{|\omega_i|\le N}\frac{1}{\langle\omega_1\rangle^2\langle\omega_2\rangle^2}\cdot\frac{1}{\langle\omega_1\rangle^2+\langle\omega_2\rangle^2+\langle\omega_1+\omega_2\rangle^2}
\gtrsim\sum_{|\omega_1|\le|\omega_2|\le N}\frac{1}{(1+|\omega_1|^2)(1+|\omega_2|^4)}.
\]
Bounding the sum below by an integral gives
$\int_0^N\frac{r}{1+r^2}\,dr\simeq\log N$.
\end{proof}

The following standard facts are recorded for reference.

\begin{lemma}\label{lem:std}
$(i)$ For any polynomial $P$, $\,\Lin_N P(\Lin_N)$ converges a.s.\ in
$\C^{-1/2-\kappa}$. $(ii)$ The expressions ${:}\Lin_N^2{:}\,\Iqc_N-3c_{N,2}\Iqc_N$
and ${:}\Lin_N^2{:}\,v_N+3c_{N,2}(v_N-\Iqc_N)$ both converge a.s.\ in
$\C^{-1-2\kappa}$ as $N\to\infty$.
\end{lemma}

\section*{Step 1: $\mu(B_\gamma)=1$}

Fix $\gamma>\tfrac12$. Let $u$ be as in \eqref{eq:u} and write $u_N=\PN u$.
Expanding the Wick cube and using $u_N=\Lin_N-\Iqc_N+v_N$,
\begin{align*}
(\log N)^{-\gamma}\Hp_3(u_N;c_N)
&=(\log N)^{-\gamma}{:}\Lin_N^3{:}
-3(\log N)^{-\gamma}{:}\Lin_N^2{:}\,\Iqc_N\\
&\quad+3(\log N)^{-\gamma}{:}\Lin_N^2{:}\,v_N
+(\log N)^{-\gamma}R_N,
\end{align*}
where $R_N$ collects the remaining mixed terms. By Lemma~\ref{lem:cube} the
first term $\to0$; by Lemma~\ref{lem:std}(i) and standard product estimates the
terms in $R_N$ $\to0$. It remains to treat
$-{:}\Lin_N^2{:}\,\Iqc_N+{:}\Lin_N^2{:}\,v_N+3c_{N,2}u_N$, which we regroup as
\[
\Bigl({:}\Lin_N^2{:}\,v_N+3c_{N,2}(v_N-\Iqc_N)\Bigr)-\Bigl({:}\Lin_N^2{:}\,\Iqc_N-3c_{N,2}\Iqc_N\Bigr).
\]
By Lemma~\ref{lem:std}(ii) this converges a.s.\ in $\C^{-1-2\kappa}$; multiplied
by $(\log N)^{-\gamma}\to0$ it tends to $0$. Pairing with $\psi$,
\[
\bigl\langle(\log N)^{-\gamma}\bigl(\Hp_3(u_N;c_N)+9c_{N,2}u_N\bigr),\psi\bigr\rangle\xrightarrow[N\to\infty]{}0\quad\text{a.s.},
\]
i.e.\ $u(t)\in B_\gamma$ for fixed $t$. Hence $\mu(B_\gamma)=1$.

\section*{Step 2: $(T_\psi^*\mu)(B_\gamma)=0$}

We must show $u+\psi\notin B_\gamma$. Expanding with $\psi_N=\PN\psi$,
\begin{align*}
&(\log N)^{-\gamma}\Hp_3(u_N+\psi_N;c_N)+(\log N)^{-\gamma}9c_{N,2}(u_N+\psi_N)\\
&\quad=\Bigl[(\log N)^{-\gamma}\Hp_3(u_N;c_N)+9(\log N)^{-\gamma}c_{N,2}u_N\Bigr]\\
&\qquad+3(\log N)^{-\gamma}{:}u_N^2{:}\,\psi_N+3(\log N)^{-\gamma}u_N\psi_N^2\\
&\qquad+(\log N)^{-\gamma}\psi_N^3+9(\log N)^{-\gamma}c_{N,2}\psi_N.
\end{align*}
The bracketed term tends to $0$ a.s.\ (Step 1). Since ${:}u_N^2{:}$ and $u_N$
converge to finite distributional limits (Lemma~\ref{lem:std}(i)), the next three
terms tend to $0$ after multiplication by $(\log N)^{-\gamma}$. The decisive term
is the last: by Lemma~\ref{lem:logdiv},
\[
9(\log N)^{-\gamma}c_{N,2}\,\langle\psi_N,\psi\rangle\ \gtrsim\ (\log N)^{-\gamma+1}\,\langle\psi_N,\psi\rangle .
\]
As $N\to\infty$, $\langle\psi_N,\psi\rangle\to\|\psi\|_{L^2}^2>0$ (because
$\psi\not\equiv0$), and $(\log N)^{-\gamma+1}\to\infty$ for $\gamma<1$. Choosing
$\tfrac12<\gamma<1$, the displayed quantity \emph{diverges}, so the full pairing
does not converge to $0$; that is, $u+\psi\notin B_\gamma$ a.s. Therefore
$(T_\psi^*\mu)(B_\gamma)=\mu(\{u:u+\psi\in B_\gamma\})=0$.

\section*{Conclusion}

By Steps 1 and 2, $\mu(B_\gamma)=1$ while $(T_\psi^*\mu)(B_\gamma)=0$. Thus $\mu$
and $T_\psi^*\mu$ are concentrated on disjoint Borel sets and are mutually
singular. The only structural inputs were that $\psi\not\equiv0$ (used solely to
ensure $\|\psi\|_{L^2}^2>0$) and that the coupling is nonzero, so the conclusion
holds for all such $\psi$. \qed

\section*{Remarks}

\begin{enumerate}\setlength\itemsep{4pt}
\item \textbf{Where dimension three enters.} The entire effect is carried by
$c_{N,2}=\E[\Lsq_N\Iqc_N]\gtrsim\log N$ (Lemma~\ref{lem:logdiv}). This logarithmic
divergence is absent in $d=2$, where $\mu$ \emph{is} quasi-invariant under smooth
shifts and even equivalent to the free field. The borderline dimension for
shift-quasi-invariance is exactly $3$, and the question is on which side it falls;
the answer is the singular side.

\item \textbf{Why the heuristic ``density'' argument fails.} Formally
$dT_\psi^*\mu/d\mu$ involves a term proportional to the would-be field
$\Lin^3-C\Lin$, which does \emph{not} define a random distribution in $d=3$ (its
covariance behaves like $|x-y|^{-3}$, non-integrable). The log-divergent constant
in the formal density is precisely what makes the candidate density blow up. The
rigorous proof above replaces the nonexistent density by the cut-off functional
defining $B_\gamma$ and extracts the divergence cleanly.

\item \textbf{Role of $\gamma$.} The window $\tfrac12<\gamma<1$ is forced:
$\gamma>\tfrac12$ kills the cube and the chaos remainders (Lemma~\ref{lem:cube}),
while $\gamma<1$ lets the shift term $(\log N)^{-\gamma}c_{N,2}\asymp(\log
N)^{1-\gamma}$ diverge. Any $\gamma$ in this range distinguishes the measures.
\end{enumerate}

\begin{thebibliography}{9}
\bibitem{HP} M.~Hairer, J.~Peroni.
  \emph{Quasi shift invariance of $\Phi^4$ measures}. To appear.
\bibitem{HKN24} M.~Hairer, S.~Kusuoka, H.~Nagoji.
  \emph{Singularity of solutions to singular SPDEs}.
  Preprint, arXiv:2409.10037, 2024.
\bibitem{Glimm68} J.~Glimm.
  \emph{Boson fields with the $\Phi^4$ interaction in three dimensions}.
  Comm.\ Math.\ Phys.\ \textbf{10} (1968), 1--47.
\bibitem{HM18} M.~Hairer, K.~Matetski.
  \emph{Discretisations of rough stochastic PDEs}.
  Ann.\ Probab.\ \textbf{46} (2018), 1651--1709.
\bibitem{EW24} S.~Esquivel, H.~Weber.
  \emph{A priori bounds for the dynamic fractional $\Phi^4$ model on $\T^3$}.
  Preprint, arXiv:2411.16536, 2024.
\bibitem{CC18} R.~Catellier, K.~Chouk.
  \emph{Paracontrolled distributions and the $3$-dimensional stochastic
  quantization equation}. Ann.\ Probab.\ \textbf{46} (2018), 2621--2679.
\bibitem{Hairer14} M.~Hairer.
  \emph{A theory of regularity structures}. Invent.\ Math.\ \textbf{198} (2014),
  269--504.
\end{thebibliography}

\end{document}